%% toc: P2 15 15.5,15.6,15.7
\section{Horn-SAT, 2-SAT and propositional decision complexity}
\label{sec:p02-w02-p02-019-sat-procedures}

Section~\ref{sec:p02-prop-decision-complexity} states which propositional
decision problems are decidable, which are \(\mathsf{NP}\)- or
\(\mathsf{coNP}\)-complete, and which restricted clause classes are decided in
linear time.  Several of those statements carry only a citation.  This section
supplies the procedures and the proofs behind them.  It writes out the forward
closure of a Horn clause list with its counter-and-queue implementation, the
implication graph of a 2-CNF list with a linear-time computation of its
strongly connected components and of a satisfying valuation, elimination for
parity systems over the two-element field, the truth-table
decision procedure, the certificate check that places \(\mathsf{SAT}\) in
\(\mathsf{NP}\), and the derivation of the \(\mathsf{coNP}\)-completeness of
\(\mathsf{TAUT}\) from the \(\mathsf{NP}\)-completeness of \(\mathsf{SAT}\).

\emph{Inputs.} Finite clause lists in the sense of
Convention~\ref{def:p02-w02-p02-019-clause-input}; finite lists of formulas of
\(\mathsf{Fm}_{\mathsf{CPL}}\); and bit strings under the encoding of
Definition~\ref{def:p02-sat-taut}.
\emph{Conclusions.} For Horn lists: the least forward closure, the exact
contradiction condition, and a decision procedure that returns either
``unsatisfiable'' or the least model, in \(O(N)\) steps
(Theorems~\ref{thm:p02-w02-p02-019-horn-closure}
and~\ref{thm:p02-w02-p02-019-horn-procedure}).  For 2-CNF lists: the
strongly connected components in linear time
(Theorem~\ref{thm:p02-w02-p02-019-scc}), the opposite-pair test and the
extraction of a satisfying valuation, again in \(O(N)\) steps
(Theorem~\ref{thm:p02-w02-p02-019-two-sat}).  For parity systems: the
correctness of Gauss--Jordan elimination over the two-element field and its
bound in field operations (Theorem~\ref{thm:p02-w02-p02-019-xor}).  For arbitrary formulas: a
truth-table bound for satisfiability, validity, equivalence and finite
consequence (Theorem~\ref{thm:p02-w02-p02-019-truth-table}),
\(\mathsf{SAT}\in\mathsf{NP}\)
(Theorem~\ref{thm:p02-w02-p02-019-sat-in-np}), and the
\(\mathsf{coNP}\)-completeness of \(\mathsf{TAUT}\)
(Theorem~\ref{thm:p02-w02-p02-019-taut-conp}).  For first-order clause sets:
refutational completeness of resolution, stated with its cited standing
(Theorem~\ref{thm:p02-w02-p02-019-fol-resolution}), and the contrast between
its one-sided refutation search and the propositional decision procedures.
The proofs of the proved results are in the annex.  They are derived from the inherited definitions
of \cref{sec:prop-syntax-semantics-01,sec:prop-syntax-semantics-02,sec:prop-syntax-semantics-04}
and from Convention~\ref{def:p02-complexity-standing}.  They are not presented
as reconstructions of the algorithms in the works cited at
Theorems~\ref{thm:p02-horn-sat} and~\ref{thm:p02-two-sat}.

\subsection{Clause inputs and the cost measure}
\label{sec:p02-w02-p02-019-inputs}

\begin{hbconvention}{def:p02-w02-p02-019-clause-input}{Clause lists, their size, and the machine model}
A \emph{clause list} is a finite sequence \(S=(C_1,\ldots,C_m)\), with
\(m\geq 0\), of clauses in the sense of
\cref{sec:prop-syntax-semantics-04}: each \(C_j\) is a finite set of literals,
presented as a list of its members without repetition.  A valuation \(v\)
satisfies \(S\) iff it satisfies every \(C_j\); by
Lemma~\ref{lem:prop-clause-satisfaction} this holds iff every \(C_j\) contains
a literal true under \(v\).  The empty list (\(m=0\)) is satisfied by every
valuation, and the empty clause by none.
\begin{enumerate}
  \item \(A_S\) is the set of variables occurring in \(S\), and \(n=|A_S|\).
        The variables of \(A_S\) are assumed renamed, in order of first
        occurrence, as \(p_1,\ldots,p_n\).  Renaming is a bijection between
        variables, so it carries satisfying valuations to satisfying
        valuations in both directions and preserves satisfiability.
  \item \(N=|C_1|+\cdots+|C_m|\) is the number of literal occurrences.  The
        size of \(S\) is \(m+N\).  If no clause is empty, then \(m\leq N\);
        and always \(n\leq N\), because every variable of \(A_S\) occurs.
        Hence, after an empty clause has been excluded, \(O(m+n+N)=O(N)\).
  \item The literal \(p_q\) is stored as the number \(2q\) and \(\neg p_q\) as
        \(2q+1\), the code of \cref{sec:prop-syntax-semantics-04}.  The
        complement \(\overline L\) of a literal is then obtained by changing
        the last binary digit.
  \item Running times are counted on a random-access machine: reading or
        writing an array cell, comparing or adding numbers of size at most
        \(2(m+n+N)+2\), and appending to or removing from the end or front of
        a list each cost one step.  An array of length \(\ell\) is created and
        initialized in \(O(\ell)\) steps.  This is the machine on which
        \textcite[121]{AspvallPlassTarjan1979} state their linear bound.
\end{enumerate}
The measure \(N\) is the one of Definition~\ref{def:p02-tractable-classes}.
The convention is needed because a linear bound is a claim about a particular
input presentation and machine, and the proofs below count steps on exactly
this one.
\index{clause list}\index{literal occurrence}
\end{hbconvention}

\subsection{Horn clauses and forward closure}
\label{sec:p02-w02-p02-019-horn}

\begin{hbdefinition}{def:p02-w02-p02-019-horn-parts}{Parts of a Horn clause}
Let \(C\) be a Horn clause (Definition~\ref{def:p02-tractable-classes}).  Its
\emph{body} is \(B(C)=\{q: \neg q\in C\}\).  If \(C\) contains a positive
literal, that literal is unique and is the \emph{head} \(\operatorname{hd}(C)\),
and \(C\) is \emph{definite}; otherwise \(C\) is \emph{negative}.  Thus
\begin{enumerate}
  \item a definite clause with \(B(C)=\{q_1,\ldots,q_k\}\) and head \(r\)
        is satisfied by \(v\) iff \(v(r)=\mathbf{T}\) or some \(v(q_i)=\mathbf{F}\),
        that is, iff it reads as the implication from all \(q_i\) to \(r\);
  \item a \emph{positive unit clause} \(\{r\}\) is the definite clause with
        empty body; it is the only kind of clause that makes an atom true
        without any premise, and it is the seed of forward chaining;
  \item a negative clause with body \(\{q_1,\ldots,q_k\}\) is satisfied by
        \(v\) iff some \(v(q_i)=\mathbf{F}\); a \emph{negative unit clause}
        \(\{\neg q\}\) has body \(\{q\}\); and the empty clause is the
        negative clause with empty body.
\end{enumerate}
A set \(X\subseteq A_S\) is \emph{closed under} a Horn list \(S\) iff, for
every definite clause \(C\) of \(S\) with \(B(C)\subseteq X\),
\(\operatorname{hd}(C)\in X\).  For any set \(X\) of variables, \(v_X\) is the
valuation with \(v_X(p)=\mathbf{T}\) iff \(p\in X\).
\index{Horn clause!definite}\index{Horn clause!negative}
\end{hbdefinition}

\begin{hbdefinition}{def:p02-w02-p02-019-forward-closure}{Least forward closure}
For a Horn list \(S\), \(\operatorname{Fwd}(S)\) is the intersection of all
subsets of \(A_S\) that are closed under \(S\).  The family is nonempty because
\(A_S\) itself is closed: every head of a clause of \(S\) occurs in \(S\).  The
intersection of closed sets is closed, so \(\operatorname{Fwd}(S)\) is the
least closed subset of \(A_S\)
(Theorem~\ref{thm:p02-w02-p02-019-horn-closure}, clause 1).  Equivalently,
it is the set of atoms obtained by starting from the heads of the positive
unit clauses and repeatedly adding the head of any definite clause whose whole
body has already been obtained; Theorem~\ref{thm:p02-w02-p02-019-horn-procedure}
proves that the procedure below computes exactly this set.
\index{forward closure}
\end{hbdefinition}

\begin{hbtheorem}{thm:p02-w02-p02-019-horn-closure}{proved=annex:p02-w02-p02-019-hc}{Forward closure and contradiction}
Let \(S\) be a Horn list and \(F=\operatorname{Fwd}(S)\).
\begin{enumerate}
  \item \(F\) is closed under \(S\) and is contained in every subset of
        \(A_S\) closed under \(S\).
  \item Every valuation satisfying \(S\) makes every atom of \(F\) true.
  \item \(v_F\) satisfies every definite clause of \(S\), and it satisfies a
        negative clause \(C\) of \(S\) iff \(B(C)\not\subseteq F\).
  \item \emph{Contradiction condition.}  \(S\) is unsatisfiable iff \(S\)
        contains a negative clause \(C\) with \(B(C)\subseteq F\), that is,
        either the empty clause or a purely negative clause all of whose
        atoms lie in \(F\).  If \(S\) is satisfiable, \(v_F\) satisfies \(S\)
        and every valuation satisfying \(S\) makes true every atom that
        \(v_F\) makes true; \(v_F\) is the least model.
  \item If \(S\) contains no positive unit clause, then \(F=\varnothing\);
        if it then also contains no empty clause, the valuation making every
        variable false satisfies \(S\).  The empty list is satisfied by every
        valuation, and \(\operatorname{Fwd}\) of it is \(\varnothing\).
\end{enumerate}
\end{hbtheorem}

\begin{hbmethod}{def:p02-w02-p02-019-horn-method}{Forward chaining with clause counters}
Input: a Horn list \(S=(C_1,\ldots,C_m)\) over \(p_1,\ldots,p_n\).
\begin{enumerate}
  \item[(0)] \emph{Empty-clause pre-check.}  For \(j=1,\ldots,m\): if
        \(|C_j|=0\), output ``unsatisfiable'' and halt.
  \item[(1)] \emph{Tables.}  Create arrays \(\mathit{mark}[1..n]\) (all
        \(0\)), \(\mathit{occ}[1..n]\) (all empty lists),
        \(\mathit{cnt}[1..m]\) (all \(0\)) and \(\mathit{hd}[1..m]\) (all
        \(0\)).  For each \(j\) and each literal of \(C_j\): if it is
        \(\neg p_q\), append \(j\) to \(\mathit{occ}[q]\) and add \(1\) to
        \(\mathit{cnt}[j]\); if it is \(p_q\), set \(\mathit{hd}[j]:=q\).
  \item[(2)] \emph{Seeds.}  Start with an empty queue \(Q\).  For each \(j\)
        with \(\mathit{cnt}[j]=0\): put \(q:=\mathit{hd}[j]\); if
        \(\mathit{mark}[q]=0\), set \(\mathit{mark}[q]:=1\) and append \(q\)
        to \(Q\).
  \item[(3)] \emph{Propagation.}  While \(Q\) is nonempty: remove the front
        atom \(q\) of \(Q\); for each \(j\) in \(\mathit{occ}[q]\), subtract
        \(1\) from \(\mathit{cnt}[j]\), and if now \(\mathit{cnt}[j]=0\),
        \(r:=\mathit{hd}[j]\neq 0\) and \(\mathit{mark}[r]=0\), set
        \(\mathit{mark}[r]:=1\) and append \(r\) to \(Q\).
  \item[(4)] \emph{Goal check.}  For each \(j\) with \(\mathit{hd}[j]=0\): if
        \(\mathit{cnt}[j]=0\), output ``unsatisfiable'' and halt.  Otherwise
        output ``satisfiable'' and the valuation that makes \(p_q\) true iff
        \(\mathit{mark}[q]=1\) and every other variable false.
\end{enumerate}
In step~(2), a clause with \(\mathit{cnt}[j]=0\) has empty body; since step~(0)
excluded the empty clause, it is a positive unit clause and
\(\mathit{hd}[j]\neq 0\).  The counter \(\mathit{cnt}[j]\) records how many
body atoms of \(C_j\) have not yet been taken from the queue.  This is forward
chaining from the positive unit clauses; read as unit propagation, removing a
true atom \(q\) shortens every clause containing \(\neg q\), and a definite
clause whose negative literals are all removed has become the unit clause of
its head.
\index{forward chaining}\index{Horn-SAT}
\end{hbmethod}

\begin{hbtheorem}{thm:p02-w02-p02-019-horn-procedure}{proved=annex:p02-w02-p02-019-hp}{Forward chaining is correct and linear}
Method~\ref{def:p02-w02-p02-019-horn-method} halts on every Horn list \(S\).  If
\(S\) contains the empty clause it outputs ``unsatisfiable'' in step~(0).
Otherwise, when step~(3) ends, the set of marked atoms is exactly
\(\operatorname{Fwd}(S)\); the method outputs ``unsatisfiable'' iff \(S\) is
unsatisfiable; and when it outputs ``satisfiable'' the valuation it returns is
\(v_{\operatorname{Fwd}(S)}\), which satisfies \(S\).  Each atom enters the
queue at most once, each occurrence of a negative literal \(\neg p_q\) in
\(C_j\) causes at most one decrement of \(\mathit{cnt}[j]\), and each positive
literal occurrence is read once.  The method runs in \(O(m+n+N)\) steps; on
lists without an empty clause this is \(O(N)\), and when an empty clause is
present step~(0) halts after at most \(N+1\) clause inspections.
\end{hbtheorem}

Theorem~\ref{thm:p02-w02-p02-019-horn-procedure} is a complete proof of the
statement of Theorem~\ref{thm:p02-horn-sat}, whose status there is a
citation.  The dual-Horn reduction of Corollary~\ref{thm:p02-dual-horn-sat}
complements every variable and then applies a linear-time Horn algorithm; that
algorithm is Method~\ref{def:p02-w02-p02-019-horn-method}, and its bound is the
one proved here.  Complementation rewrites each literal occurrence once, so the
dual-Horn bound is also \(O(N)\).  The least closure proved for the circuit
instance in Annex~\ref{annex:p14.dc.horn} is the special case of clauses~1, 2
and~3 of Theorem~\ref{thm:p02-w02-p02-019-horn-closure} for that instance; the
\(\mathsf P\)-completeness of Horn satisfiability is
Theorem~\ref{thm:p14-dc-horn}.  Read as a propositional definite program, with
negative clauses as goals, a Horn list has
\(\operatorname{Fwd}(S)=M_{P_S}\), the least fixed point of the
immediate-consequence operator \(T_{P_S}\) of Part~9
(\cref{sec:p09-logic-programming}, Definition~\ref{def:p09-definite-program});
this is Theorem~\ref{thm:p02-w02b-p02-015-horn-lfp}, which also proves the
variable-free case of the least-model Theorem~\ref{thm:p09-definite-fixpoint}.

\begin{hbexample}{def:p02-w02-p02-019-horn-example}{A satisfiable Horn list and a contradictory extension}
Let \(S\) be the list
\[
 C_1=\{p_1\},\quad
 C_2=\{\neg p_1,p_2\},\quad
 C_3=\{\neg p_2,\neg p_4,p_3\},\quad
 C_4=\{\neg p_1,\neg p_3\}.
\]
Here \(m=4\), \(n=4\) and \(N=1+2+3+2=8\).  Step~(1) gives
\(\mathit{cnt}=(0,1,2,2)\), \(\mathit{hd}=(1,2,3,0)\),
\(\mathit{occ}[1]=(2,4)\), \(\mathit{occ}[2]=(3)\), \(\mathit{occ}[3]=(4)\)
and \(\mathit{occ}[4]=(3)\).  Step~(2) marks \(p_1\), the head of the
positive unit clause \(C_1\).  Removing \(p_1\) lowers \(\mathit{cnt}[2]\) to
\(0\), which marks \(p_2\), and lowers \(\mathit{cnt}[4]\) to \(1\).
Removing \(p_2\) lowers \(\mathit{cnt}[3]\) to \(1\); the queue is then empty.
So \(\operatorname{Fwd}(S)=\{p_1,p_2\}\).  The only negative clause is
\(C_4\), and \(\mathit{cnt}[4]=1\), because its body atom \(p_3\) is not in
the closure.  The output is ``satisfiable'' with
\(p_1,p_2\) true and \(p_3,p_4\) false; this valuation satisfies all four
clauses.

Append the positive unit clause \(C_5=\{p_4\}\).  The seeds are now \(p_1\)
and \(p_4\).  Removing \(p_1\) marks \(p_2\) and sets \(\mathit{cnt}[4]=1\);
removing \(p_4\) sets \(\mathit{cnt}[3]=1\); removing \(p_2\) sets
\(\mathit{cnt}[3]=0\) and marks \(p_3\); removing \(p_3\) sets
\(\mathit{cnt}[4]=0\).  Now \(\operatorname{Fwd}=\{p_1,p_2,p_3,p_4\}\)
contains the whole body \(\{p_1,p_3\}\) of the negative clause \(C_4\), and
step~(4) outputs ``unsatisfiable''.  By clause~2 of
Theorem~\ref{thm:p02-w02-p02-019-horn-closure} every model would make \(p_1\)
and \(p_3\) true and so falsify \(C_4\).
\end{hbexample}

\begin{hbwarning}{def:p02-w02-p02-019-horn-limit}{Where forward closure stops working}
The list consisting of the single clause \(\{p_1,p_2\}\) is not Horn.  It is
satisfiable, but its models \(v_{\{p_1\}}\) and \(v_{\{p_2\}}\) have no
common part that is again a model: the valuation making both atoms false
falsifies it.  So the list has no least model, and no procedure that computes
a least closed set and tests it can be correct beyond the Horn class.  Forward
chaining applied to it would start from no seed and return the all-false
valuation, which is not a model.
The linear bound also depends on the input convention: renaming arbitrary
variable names or indices to \(p_1,\ldots,p_n\) is not counted in
Theorem~\ref{thm:p02-w02-p02-019-horn-procedure}.  For names compared as
strings it needs a sort or a search structure, and that cost has to be added.
\end{hbwarning}

\subsection{Implication graphs and 2-SAT}
\label{sec:p02-w02-p02-019-two-sat}

\begin{hbdefinition}{def:p02-w02-p02-019-implication-graph}{Implication graph, strong components and condensation}
Let \(S=(C_1,\ldots,C_m)\) be a 2-CNF list over \(p_1,\ldots,p_n\) with no empty
clause.  Its \emph{implication graph} \(G_S\) has the \(2n\) literals over
\(p_1,\ldots,p_n\) as vertices and the following edge list: a two-literal
clause \(\{L,K\}\) contributes the edges \(\overline L\to K\) and
\(\overline K\to L\); a unit clause \(\{L\}\) is read as \(L\lor L\) and
contributes the edge \(\overline L\to L\) twice.  Thus \(G_S\) has exactly
\(2m\) edges, counted with repetition.  The \emph{reverse graph}
\(G_S^{\mathrm R}\) has the same vertices and an edge \(K\to L\) for each edge
\(L\to K\) of \(G_S\).

For vertices \(X,Y\), write \(X\leadsto Y\) iff there is a directed path, of
length \(0\) or more, from \(X\) to \(Y\).  The relation \(X\leadsto Y\) and
\(Y\leadsto X\) is an equivalence relation; its classes are the \emph{strong
components}, and \(C(X)\) is the strong component of \(X\).  The
\emph{condensation} has the strong components as vertices and an edge from
\(C(X)\) to \(C(Y)\) whenever \(G_S\) has an edge \(X\to Y\) with
\(C(X)\neq C(Y)\).  A numbering of the components is \emph{topological} iff
every condensation edge goes from a smaller to a larger number.
\index{implication graph}\index{strongly connected component}
\end{hbdefinition}

This is the graph of Theorem~\ref{thm:p02-two-sat}; the construction with
\(2n\) vertices and \(2|C|\) edges, and the reading of a unit clause \(u\) as
\(u\lor u\), are those of \textcite[121]{AspvallPlassTarjan1979}.

\begin{hbmethod}{def:p02-w02-p02-019-scc-method}{Strong components by two depth-first passes}
Input: a finite directed graph \(G\) given by adjacency lists, and its reverse
\(G^{\mathrm R}\).
\begin{enumerate}
  \item \emph{First pass.}  Colour every vertex white and set \(t:=0\).  For
        each vertex \(u\) in a fixed order, if \(u\) is white, call
        \(\mathrm{Visit}(u)\).  Here \(\mathrm{Visit}(u)\) colours \(u\) grey;
        for each out-neighbour \(w\) of \(u\) in \(G\), in adjacency order,
        calls \(\mathrm{Visit}(w)\) if \(w\) is white; then colours \(u\)
        black, sets \(t:=t+1\) and \(f(u):=t\), and appends \(u\) to a list
        \(\mathit{fin}\).  The number \(f(u)\) is the \emph{finishing
        number} of \(u\).
  \item \emph{Second pass.}  Mark every vertex unvisited and set \(k:=0\).
        Run through \(\mathit{fin}\) from its last entry to its first, that
        is, in decreasing finishing number.  When the current vertex \(r\) is
        unvisited, set \(k:=k+1\) and call \(\mathrm{Visit}^{\mathrm R}(r)\),
        which marks a vertex visited, gives it the number
        \(\operatorname{idx}:=k\), and calls itself on every unvisited
        out-neighbour in \(G^{\mathrm R}\).  The set of vertices visited by
        one such call from \(r\) is the \emph{tree} of \(r\).
\end{enumerate}
The recursion is implemented with an explicit stack of pairs (vertex, position
in its adjacency list), so each step of a call costs \(O(1)\) on the machine of
Convention~\ref{def:p02-w02-p02-019-clause-input}.
\index{depth-first search}
\end{hbmethod}

\begin{hbtheorem}{thm:p02-w02-p02-019-scc}{proved=annex:p02-w02-p02-019-scc}{Two-pass strong components}
Let \(G\) have \(V\) vertices and \(E\) edges.  Each tree of the second pass of
Method~\ref{def:p02-w02-p02-019-scc-method} is exactly one strong component of
\(G\), and every strong component is one tree.  The number
\(\operatorname{idx}\) of the trees is a topological numbering of the
condensation: if \(G\) has an edge \(X\to Y\) and \(C(X)\neq C(Y)\), then
\(\operatorname{idx}(X)<\operatorname{idx}(Y)\).  The method runs in
\(O(V+E)\) steps.
\end{hbtheorem}

\begin{hbmethod}{def:p02-w02-p02-019-two-sat-method}{The 2-SAT test with assignment extraction}
Input: a 2-CNF list \(S\) over \(p_1,\ldots,p_n\).
\begin{enumerate}
  \item[(0)] If some clause is empty, output ``unsatisfiable'' and halt.
  \item[(1)] Build the adjacency lists of \(G_S\) and \(G_S^{\mathrm R}\) by
        one pass over the clauses, two edges per clause.
  \item[(2)] Run Method~\ref{def:p02-w02-p02-019-scc-method} on \(G_S\) to
        obtain \(\operatorname{idx}\).
  \item[(3)] \emph{Opposite-pair check.}  For \(q=1,\ldots,n\): if
        \(\operatorname{idx}(p_q)=\operatorname{idx}(\neg p_q)\), output
        ``unsatisfiable'' and halt.
  \item[(4)] \emph{Extraction.}  Output ``satisfiable'' and the valuation
        \(v_S\) with \(v_S(p_q)=\mathbf{T}\) iff
        \(\operatorname{idx}(p_q)>\operatorname{idx}(\neg p_q)\), for
        \(q=1,\ldots,n\), and every other variable false.
\end{enumerate}
A literal is thus made true when its component comes later than the
component of its complement in the topological numbering.  This is the
marking of \textcite[122]{AspvallPlassTarjan1979}, who process the components
in reverse topological order and make the first unmarked component of each
complementary pair true.
\index{2-SAT}
\end{hbmethod}

\begin{hbtheorem}{thm:p02-w02-p02-019-two-sat}{proved=annex:p02-w02-p02-019-2sat}{The 2-SAT test}
Let \(S\) be a 2-CNF list over \(p_1,\ldots,p_n\) with no empty clause.
\begin{enumerate}
  \item A valuation satisfies \(S\) iff, for every edge \(X\to Y\) of
        \(G_S\), \(v\entails X\) implies \(v\entails Y\).
  \item For every edge \(X\to Y\) of \(G_S\), \(\overline Y\to\overline X\) is
        an edge of \(G_S\).  Consequently \(C(\overline X)=\{\overline Z:
        Z\in C(X)\}\) for every vertex \(X\).
  \item If \(C(p_q)=C(\neg p_q)\) for some \(q\), then \(S\) is
        unsatisfiable.
  \item If \(C(p_q)\neq C(\neg p_q)\) for every \(q\), then the valuation
        \(v_S\) of Method~\ref{def:p02-w02-p02-019-two-sat-method}, step~(4),
        satisfies \(S\).
  \item Method~\ref{def:p02-w02-p02-019-two-sat-method} outputs
        ``unsatisfiable'' iff \(S\) is unsatisfiable, and otherwise outputs a
        satisfying valuation.  It runs in \(O(m+n)\) steps, hence in \(O(N)\)
        steps.  The empty list (\(m=n=0\)) yields ``satisfiable'' with the
        empty assignment to the empty variable set.
\end{enumerate}
\end{hbtheorem}

Clauses~3 and~4 are the criterion stated in Theorem~\ref{thm:p02-two-sat} and
proved, with an arbitrary topological rank, in Annex~\ref{annex:p04.dh.two-sat};
the proof here uses the particular rank \(\operatorname{idx}\) produced by
Method~\ref{def:p02-w02-p02-019-scc-method}, so that the satisfying valuation is
computed, not only shown to exist.  Clause~5 is a complete proof of the
\(O(N)\) bound of Theorem~\ref{thm:p02-two-sat}, which there carries the
citation \cite[thm.~1, p.~122]{AspvallPlassTarjan1979}; that note bounds its
algorithm by \(O(n+m)\) with \(m\) the number of graph edges
\cite[123]{AspvallPlassTarjan1979}.  The \(\mathsf{NL}\)-completeness of
2-CNF satisfiability is Theorem~\ref{thm:p04-dh-two-sat}.  A 2-CNF clause
with two positive literals is not a definite clause of
Definition~\ref{def:p09-definite-program}, so the least-model reading of
Part~9 (\cref{sec:p09-logic-programming}), which covers the Horn lists through
Theorem~\ref{thm:p02-w02b-p02-015-horn-lfp}, does not extend to 2-CNF lists
in general (Warning~\ref{def:p02-w02b-p02-015-two-sat-programs}).

\begin{hbexample}{def:p02-w02-p02-019-two-sat-example}{A 2-SAT run and a contradictory list}
Let \(S\) consist of \(\{p_1,p_2\}\), \(\{\neg p_1,p_2\}\) and
\(\{\neg p_2,p_3\}\).  The edges of \(G_S\) are
\[
 \neg p_1\to p_2,\quad \neg p_2\to p_1,\quad
 p_1\to p_2,\quad \neg p_2\to\neg p_1,\quad
 p_2\to p_3,\quad \neg p_3\to\neg p_2.
\]
Take the vertex order \(p_1,\neg p_1,p_2,\neg p_2,p_3,\neg p_3\) and the
adjacency order above.  The first pass visits \(p_1,p_2,p_3\) in one call and
finishes \(p_3,p_2,p_1\); then it finishes \(\neg p_1\), \(\neg p_2\) and
\(\neg p_3\).  In the second pass the decreasing finishing order is
\(\neg p_3,\neg p_2,\neg p_1,p_1,p_2,p_3\).  Each in-neighbour in \(G_S\) of
each of these vertices is already visited when the vertex is reached, so every
tree is a single vertex, and
\[
 \begin{array}{lll}
 \operatorname{idx}(\neg p_3)=1, & \operatorname{idx}(\neg p_2)=2, &
 \operatorname{idx}(\neg p_1)=3,\\
 \operatorname{idx}(p_1)=4, & \operatorname{idx}(p_2)=5, &
 \operatorname{idx}(p_3)=6.
 \end{array}
\]
Every edge listed goes from a smaller to a larger number.  No opposite pair
shares a component, and step~(4) makes \(p_1\), \(p_2\) and \(p_3\) true
(\(4>3\), \(5>2\), \(6>1\)).  This valuation satisfies the three clauses.

Add the clause \(\{\neg p_2,\neg p_3\}\), with edges \(p_2\to\neg p_3\) and
\(p_3\to\neg p_2\).  Now
\(p_2\to p_3\to\neg p_2\to p_1\to p_2\) is a cycle, so \(p_2\) and
\(\neg p_2\) lie in one strong component and step~(3) outputs
``unsatisfiable''.  Directly: the first two clauses force \(p_2\), and then
\(\{\neg p_2,p_3\}\) and \(\{\neg p_2,\neg p_3\}\) force \(p_3\) and
\(\neg p_3\).
\end{hbexample}

\begin{hbwarning}{def:p02-w02-p02-019-three-literals}{Three literals per clause}
A clause \(\{L,K,M\}\) says that the falsity of two literals forces the
third.  An edge of an implication graph has one literal as antecedent, and no
set of such implications has exactly the models of the clause
\(\{p_1,p_2,p_3\}\).  Indeed, let \(X\to Y\) be an implication between
literals over \(p_1,p_2,p_3\) with \(X\neq Y\), and let \(c\) be an atom
different from the atoms of \(X\) and \(Y\).  Making \(X\) true, \(Y\) false
when \(Y\) is over another atom, and \(c\) true gives a model of the clause in
which \(X\to Y\) fails; when \(Y=\overline X\), making \(X\) and \(c\) true
does the same.  So every implication valid in all models of the clause has
the form \(X\to X\), and such implications are also true under the all-false
valuation, which falsifies the clause.  The SCC method therefore does not
extend to 3-CNF.
Satisfiability of 3-CNF formulas is \(\mathsf{NP}\)-complete
\cite[thm.~2.10]{AroraBarak2009}, so a polynomial-time test for it would give
\(\mathsf P=\mathsf{NP}\).
\end{hbwarning}

\subsection{Parity systems and Gaussian elimination}
\label{sec:p02-w02-p02-019-xor}

\begin{hbconvention}{def:p02-w02-p02-019-parity-input}{Parity systems and their dense presentation}
An XOR instance in the sense of Definition~\ref{def:p02-tractable-classes}
is presented as \(m\geq 0\) equations in the variables \(p_1,\ldots,p_n\)
(\(n\geq 0\)) over the two-element field \(\{0,1\}\), whose addition is
addition modulo \(2\):
\[
  a_{j1}x_1+\cdots+a_{jn}x_n=b_j\qquad(1\leq j\leq m),
\]
with all \(a_{ji},b_j\in\{0,1\}\).  It is stored densely as the \(m\) rows
\((a_{j1},\ldots,a_{jn}\mid b_j)\) of an \(m\times(n+1)\) array, each row
reached through a row pointer.  A valuation \(v\) \emph{satisfies} equation
\(j\) iff, writing \(x_i=1\) when \(v(p_i)=\mathbf{T}\) and \(x_i=0\)
otherwise, the number of \(i\) with \(a_{ji}=1\) and \(x_i=1\) has the parity
of \(b_j\).  A \emph{field operation} is one addition of two entries.  Other
steps are counted on the machine of
Convention~\ref{def:p02-w02-p02-019-clause-input}.
\index{XOR-SAT!dense presentation}
\end{hbconvention}

\begin{hbmethod}{def:p02-w02-p02-019-gauss-method}{Elimination over the two-element field}
Set \(r:=1\) and the pivot list empty.  For \(c=1,\ldots,n\) in turn:
\begin{enumerate}
  \item Search rows \(r,\ldots,m\) for a row \(i\) with \(a_{ic}=1\).  If
        there is none, go to the next \(c\).
  \item Exchange the row pointers of rows \(i\) and \(r\).
  \item For every row \(i'\neq r\) with \(a_{i'c}=1\), replace row \(i'\)
        by its entrywise sum with row \(r\) (all \(n+1\) entries).
  \item Append \((r,c)\) to the pivot list and set \(r:=r+1\).
\end{enumerate}
When the loop ends, let \(t=r-1\) and let \((1,c_1),\ldots,(t,c_t)\) be the
pivot list.  If some row \(j>t\) has \(b_j=1\), answer ``unsatisfiable''.
Otherwise answer ``satisfiable'' and return the valuation \(v^{*}\) with
\(v^{*}(p_{c_s})=\mathbf{T}\) iff \(b_s=1\) for \(s\leq t\), and every other
variable false.  The entries named in the final test and in \(v^{*}\) are
those of the array after the loop.
\index{Gaussian elimination!over the two-element field}
\end{hbmethod}

\begin{hbtheorem}{thm:p02-w02-p02-019-xor}{proved=annex:p02-w02-p02-019-xor}{Elimination decides XOR-SAT}
Run Method~\ref{def:p02-w02-p02-019-gauss-method} on a parity system with
\(m\) equations and \(n\) variables.
\begin{enumerate}
  \item The system is satisfiable iff the method answers ``satisfiable'', and
        then \(v^{*}\) satisfies every equation.
  \item The number \(t\) of pivots satisfies \(t\leq\min\{m,n\}\), and the
        method performs at most \(t(m-1)(n+1)\) field operations, hence
        \(O(mn\min\{m,n\})\) field operations when \(m,n\geq 1\).
  \item All other work takes \(O(mn+m+n)\) steps, so the total running time,
        excluding the reading of the input array, is
        \(O(mn\min\{m,n\}+m+n)\).
  \item If \(m=0\), the method answers ``satisfiable''; if \(n=0\), it
        answers ``unsatisfiable'' iff some \(b_j=1\).
\end{enumerate}
\end{hbtheorem}

This is a complete proof of the bound stated in
Theorem~\ref{thm:p02-xor-sat} and in the XOR row of
Table~\ref{tab:p02-tractable-sat}, whose status there is a citation.  The
count is in field operations on the dense array, which has \(m(n+1)\) cells,
so it is not comparable with the \(O(N)\) bounds for Horn and 2-CNF lists:
those count literal occurrences of a clause list, while a parity equation on
many variables is one row here.

\begin{hbexample}{def:p02-w02-p02-019-xor-example}{An inconsistent parity system and a consistent variant}
Take \(n=3\) and the rows \((1,1,0\mid 1)\), \((0,1,1\mid 1)\),
\((1,0,1\mid 1)\), that is, \(x_1+x_2=1\), \(x_2+x_3=1\), \(x_1+x_3=1\).
For \(c=1\) the pivot is row \(1\); row \(3\) becomes \((0,1,1\mid 0)\).
For \(c=2\) the search starts at row \(2\), which is the pivot; row \(1\)
becomes \((1,0,1\mid 0)\) and row \(3\) becomes \((0,0,0\mid 1)\).  For
\(c=3\) no row from row \(3\) on has a \(1\) in column \(3\).  Now \(t=2\)
and row \(3\) has \(b_3=1\): the answer is ``unsatisfiable''.  Indeed the sum
of the three equations is \(0=1\).

Change the third right-hand side to \(0\).  For \(c=1\) row \(3\) becomes
\((0,1,1\mid 1)\); for \(c=2\) rows \(1\) and \(3\) become
\((1,0,1\mid 0)\) and \((0,0,0\mid 0)\); column \(3\) has no pivot.  The
pivots are \((1,1)\) and \((2,2)\), and no row below row \(2\) has right-hand
side \(1\).
The answer is ``satisfiable'' and \(v^{*}\) makes \(p_1\) false, \(p_2\)
true and \(p_3\) false.  Check: \(0+1=1\), \(1+0=1\), \(0+0=0\).
\end{hbexample}

\subsection{Decision procedures, SAT, TAUT and equivalence}
\label{sec:p02-w02-p02-019-complexity}

\begin{hbmethod}{def:p02-w02-p02-019-truth-table-method}{Truth-table decision procedure}
The input is a list of formulas \(\varphi_1,\ldots,\varphi_r\) of
\(\mathsf{Fm}_{\mathsf{CPL}}\), with \(r\geq 1\), whose jointly occurring variables are renamed as
\(p_1,\ldots,p_k\) (\(k\geq 0\)); \(\ell=|\varphi_1|+\cdots+|\varphi_r|\) is
their total number of symbol occurrences.  For each of the \(2^k\)
assignments \(u:\{p_1,\ldots,p_k\}\to\{\mathbf{F},\mathbf{T}\}\), enumerated by a
binary counter, compute \(\widehat{u^{\circ}}(\varphi_i)\) for every \(i\),
where \(u^{\circ}\) extends \(u\) by making every other variable false.  Each
value is computed by one left-to-right scan of \(\varphi_i\) with a value
stack and an operator stack: a variable or \(\bottom\) pushes its value, which
then absorbs every \(\neg\) waiting on top of the operator stack; a closing
parenthesis combines the top two values by the binary connective between them
and again absorbs waiting negations.  Then answer:
\begin{enumerate}
  \item \(\varphi_1\) is satisfiable iff some \(u\) gives it \(\mathbf{T}\);
  \item \(\varphi_1\) is valid iff every \(u\) gives it \(\mathbf{T}\);
  \item \(\varphi_1\equiv\varphi_2\) iff every \(u\) gives them the same value;
  \item \(\{\varphi_1,\ldots,\varphi_{r-1}\}\entails\varphi_r\) iff no \(u\)
        gives \(\mathbf{T}\) to \(\varphi_1,\ldots,\varphi_{r-1}\) and \(\mathbf{F}\)
        to \(\varphi_r\).
\end{enumerate}
\index{truth table!decision procedure}
\end{hbmethod}

\begin{hbtheorem}{thm:p02-w02-p02-019-truth-table}{proved=annex:p02-w02-p02-019-tt}{Truth tables decide}
The four answers of Method~\ref{def:p02-w02-p02-019-truth-table-method} are
correct for satisfiability, validity, logical equivalence and semantic
consequence from a finite set.  The method runs in \(O(2^k\ell)\) steps on the
machine of Convention~\ref{def:p02-w02-p02-019-clause-input}.  In particular
each of the four questions is decided by a terminating algorithm, in
\(O(2^k|\varphi|)\) steps for a single formula \(\varphi\) with \(k\)
variables.
\end{hbtheorem}

Theorem~\ref{thm:p02-w02-p02-019-truth-table} is a complete proof of
Theorem~\ref{thm:p02-decidability}, whose status there is a citation.  The
exponent is the number of distinct variables, not the length; a formula of any
length with few variables is decided quickly.

\begin{hbtheorem}{thm:p02-w02-p02-019-sat-in-np}{proved=annex:p02-w02-p02-019-np}{SAT is in NP under the fixed encoding}
Let \(x\in\{0,1\}^*\) have length \(n\).  Let \(R(x,y)\) hold iff \(x\) is the
encoding, in the sense of Definition~\ref{def:p02-sat-taut}, of a formula
\(\varphi\) whose distinct variables in order of first occurrence are
\(p_{a_1},\ldots,p_{a_k}\), \(y\) has length exactly \(k\), and \(\varphi\) is
true under the valuation \(v_y\) with \(v_y(p_{a_i})=\mathbf{T}\) iff the \(i\)-th
bit of \(y\) is \(1\), and every other variable false.  Then \(R\) is decided
by a deterministic multitape Turing machine in time polynomial in \(n\), every
\(y\) with \(R(x,y)\) has length at most \(n\), and \(x\in\mathsf{SAT}\) iff some
\(y\) of length at most \(n\) satisfies \(R(x,y)\).  Hence
\(\mathsf{SAT}\in\mathsf{NP}\) in the sense of
Convention~\ref{def:p02-complexity-standing}.  Malformed strings satisfy
\(R\) with no \(y\).
\end{hbtheorem}

\paragraph{Where the hardness half is proved.}
Theorem~\ref{thm:p02-sat-np-complete} has two halves.  Membership is
Theorem~\ref{thm:p02-w02-p02-019-sat-in-np}.  Hardness, that every language in
\(\mathsf{NP}\) reduces to \(\mathsf{SAT}\) by a polynomial-time many-one
reduction, is the Cook--Levin Theorem~\ref{thm:p04-cook-levin} of Part~4.
Its written proof under the Part~4 conventions, where \(\mathsf{NP}\) is
defined by nondeterministic machines with read-only input, is
Annex~\ref{annex:p02-w02b-p02-015-cook-levin}.  That proof transfers a
Part~4 language to the certificate form of
Convention~\ref{def:p02-complexity-standing}
(Lemma~\ref{thm:p02-w02b-p02-015-bridge}) and then applies the tableau
formula of Theorem~\ref{thm:p02-w02-p02-018-cook-levin}, which for every
\(x\) yields in polynomial time an encoded formula \(\varphi_x\) with
\(x\in L\) iff \(\varphi_x\) is satisfiable.  The cited formulation
\cite[thm.~2.10]{AroraBarak2009} takes \(\mathsf{SAT}\) to be the
satisfiable CNF formulas; the language here contains all satisfiable encoded
formulas, and membership for all formulas is the theorem above.  Part~9
presents the same tableau method in the form of Cook's paper
(\cref{sec:p09-presentation-n-c-026}).

\begin{hblemma}{thm:p02-w02-p02-019-reductions}{proved=annex:p02-w02-p02-019-red}{Complements and composites of reductions}
Let \(A,B,D\subseteq\{0,1\}^*\), and write \(\overline A=\{0,1\}^*\setminus A\).
\begin{enumerate}
  \item If \(f\) is a polynomial-time many-one reduction from \(A\) to \(B\),
        then the same \(f\) is one from \(\overline A\) to \(\overline B\).
  \item If \(f\) reduces \(A\) to \(B\) and \(g\) reduces \(B\) to \(D\), both
        in polynomial time, then \(g\circ f\) reduces \(A\) to \(D\) in
        polynomial time.
  \item If \(A\) reduces to \(B\) in polynomial time and
        \(B\in\mathsf{NP}\), then \(A\in\mathsf{NP}\).
\end{enumerate}
\end{hblemma}

\begin{hblemma}{thm:p02-w02-p02-019-negation-maps}{proved=annex:p02-w02-p02-019-neg}{Negation maps}
Let \(\nu\) map the encoding of a formula \(\varphi\) to the encoding of
\(\neg\varphi\), obtained by writing the block of \(\neg\) in front.  Define
\(g_1(x)=\nu(x)\) and \(g_2(x)=\nu(x)\) when \(x\) encodes a formula; for a
malformed \(x\), let \(g_1(x)\) be the encoding of \(p_0\) and \(g_2(x)\) the
encoding of \(\neg\bottom\).  Then \(\varphi\) is valid iff \(\neg\varphi\) is
unsatisfiable; \(g_1\) reduces \(\mathsf{TAUT}\) to
\(\overline{\mathsf{SAT}}\); \(g_2\) reduces \(\overline{\mathsf{SAT}}\) to
\(\mathsf{TAUT}\); and both maps are computable in polynomial time.
\end{hblemma}

\begin{hbtheorem}{thm:p02-w02-p02-019-taut-conp}{proved=annex:p02-w02-p02-019-taut}{TAUT is coNP-complete}
Assume Theorem~\ref{thm:p02-sat-np-complete}.  Then \(\mathsf{TAUT}\in\mathsf{coNP}\),
and every language in \(\mathsf{coNP}\) has a polynomial-time many-one
reduction to \(\mathsf{TAUT}\).  Thus \(\mathsf{TAUT}\) is
\(\mathsf{coNP}\)-complete, which is
Corollary~\ref{thm:p02-taut-conp-complete}.
\end{hbtheorem}

The membership half uses only
Theorem~\ref{thm:p02-w02-p02-019-sat-in-np}.  The hardness half uses the
hardness of \(\mathsf{SAT}\), proved at the place named above.  The same
negation step is \cite[ex.~2.23]{AroraBarak2009}, with \(\mathsf{coNP}\) as in
\cite[def.~2.21, def.~2.22]{AroraBarak2009}.

\begin{hbwarning}{def:p02-w02-p02-019-sat-vs-equivalence}{Satisfiability, validity and equivalence are different questions}
\(\mathsf{SAT}\) asks whether \emph{some} valuation makes one formula true.
Validity asks whether \emph{every} valuation does.  Logical equivalence of two
arbitrary formulas, \(\varphi\equiv\psi\) in the sense of
\cref{sec:prop-syntax-semantics-02}, asks whether \emph{every} valuation gives
them the same value; it holds iff
\((\varphi\to\psi)\land(\psi\to\varphi)\), the expansion of
\(\varphi\leftrightarrow\psi\), is valid.  So equivalence is a validity
question about a derived formula: its certificates are certificates of
\emph{non}-equivalence, namely one valuation on which the formulas differ, and
\(\varphi\not\equiv\psi\) iff \(\neg(\varphi\leftrightarrow\psi)\) is
satisfiable.  Equivalence is \(\mathsf{coNP}\)-complete
(Theorem~\ref{thm:p14-dc-finite-cpl}), and
Corollary~\ref{thm:p02-w02-p02-019-np-conp} shows that it lies in
\(\mathsf{NP}\) only if \(\mathsf{NP}=\mathsf{coNP}\).
Two satisfiable formulas need not be equivalent: \(p_0\) and \(p_1\) are both
satisfiable, and the valuation making \(p_0\) true and \(p_1\) false
separates them.  Transformations that preserve satisfiability, such as the
definitional CNF of Theorem~\ref{thm:p02-tseitin}, are therefore not
equivalence transformations, and a SAT answer about a transformed formula is
not an answer about equivalence with the original.
\end{hbwarning}

\begin{hbcorollary}{thm:p02-w02-p02-019-np-conp}{proved=annex:p02-w02-p02-019-npco}{Validity, equivalence and NP}
Assume Theorem~\ref{thm:p02-sat-np-complete}.  If \(\mathsf{TAUT}\in\mathsf{NP}\),
then \(\mathsf{NP}=\mathsf{coNP}\); and if the equivalence language of
Theorem~\ref{thm:p14-dc-finite-cpl} is in \(\mathsf{NP}\), then
\(\mathsf{NP}=\mathsf{coNP}\).  Conversely, if \(\mathsf{NP}=\mathsf{coNP}\),
both languages are in \(\mathsf{NP}\).
\end{hbcorollary}

The classification leaves open whether \(\mathsf P=\mathsf{NP}\) and whether
\(\mathsf{NP}=\mathsf{coNP}\).  Part~4 places these questions in its class
ladder (\cref{def:p04-complexity-convention,thm:p04-cook-levin}).  Part~12
treats satisfiability procedures as algorithms with certificates: the DPLL
core and its unit propagation are in \cref{sec:p12-sat-and-certificates}
(Definition~\ref{def:p12-dpll-state}, Method~\ref{def:p12-dpll},
Theorem~\ref{thm:p12-dpll-correctness}), and the standing of SAT and UNSAT
evidence is Definition~\ref{def:p12-sat-certificate-boundary}.  Part~9's
complexity presentation (\cref{sec:p09-presentation-n-c-026}) also
distinguishes Cook's original query reduction from the many-one reduction used
here.

\subsection{Refutational completeness and the first-order boundary}
\label{sec:p02-w02-p02-019-first-order-boundary}

Propositional resolution is refutationally complete for finite clause sets
(Theorem~\ref{thm:prop-resolution-complete}, proved in
\cref{sec:prop-metatheory-04}).  For a finite set \(S\) of clauses over
\(k\) variables, every clause that resolution can derive from \(S\) is a set
of literals over those same \(k\) variables, so at most \(2^{2k}\) clauses can
ever be derived and saturation by resolution halts.  Resolution is therefore a
decision procedure for propositional unsatisfiability, but the count of
candidate clauses is exponential in \(k\), and nothing proved here shows that
any procedure decides the question in polynomial time; that is the question
\(\mathsf P=\mathsf{NP}\) left open above.  The first-order counterpart keeps
completeness and loses the finite search space.

\begin{hbtheorem}{thm:p02-w02-p02-019-fol-resolution}{cited=Robinson1965}{Refutational completeness of first-order resolution}
Let \(S\) be a finite equality-free set of first-order clauses.  After adding
a constant if the signature has no ground term, \(S\) is not satisfiable iff
resolution with factoring derives the empty clause.
\end{hbtheorem}

This is Theorem~\ref{thm:fol-resolution} of
\cref{sec:fol-cardinality-normalforms-04}, with the same hypotheses and the
same rules (first-order Resolution and Factoring on renamed-apart parents,
with most general unifiers), stated here with its standing.  The standing is
a citation \cite{Robinson1965,BachmairGanzinger2001}.  The argument written
after Theorem~\ref{thm:fol-resolution} is a proof sketch: the soundness
theorem for first-order resolution and factoring stated just before it gives
the right-to-left direction; for the left-to-right direction it combines the
Herbrand theorem for clauses (Theorem~\ref{thm:herbrand-clauses}), the
propositional Theorem~\ref{thm:prop-resolution-complete} and the lifting
lemma (Theorem~\ref{thm:resolution-lifting}).  The Herbrand theorem and the
lifting lemma are themselves cited there, so the first-order result is not
proved in this book from proved parts alone.  The exclusion of equality is a
hypothesis, not a simplification: bare resolution is incomplete once logical
equality occurs, as the paragraph following Theorem~\ref{thm:fol-resolution}
records.

The theorem gives a one-sided test.  Together with soundness, it says that a
finite equality-free clause set is unsatisfiable exactly when some finite
derivation reaches the empty clause; it gives no bound on the length of that
derivation, and the clauses derivable from \(S\) need not form a finite set.
For example, from the satisfiable set consisting of \(\{P(a)\}\) and
\(\{\neg P(x),P(f(x))\}\), resolution derives \(\{P(f(a))\}\), then
\(\{P(f(f(a)))\}\), and so on without end, so saturation does not halt.
For first-order sentences
the decision problem itself has a negative solution: validity is not
decidable (Theorem~\ref{thm:fol-validity-undecidable}, cited to Church and
Turing), and the set of valid sentences is computably enumerable but not
decidable (Corollary~\ref{thm:fol-validity-enumerable}, with its proof from
completeness of the first-order calculus).  So the three boundaries of this
section are of different kinds: Horn and 2-CNF satisfiability are decided in
linear time
(Theorems~\ref{thm:p02-w02-p02-019-horn-procedure}
and~\ref{thm:p02-w02-p02-019-two-sat}), and parity systems in polynomial time
(Theorem~\ref{thm:p02-w02-p02-019-xor}); unrestricted propositional
satisfiability is decided, and is \(\mathsf{NP}\)-complete
(Theorem~\ref{thm:p02-sat-np-complete}); first-order unsatisfiability of
equality-free clause sets has a complete refutation search
(Theorem~\ref{thm:p02-w02-p02-019-fol-resolution}) while first-order validity
has no decision procedure at all.
